\documentclass[10pt,a4paper]{amsart}
\usepackage[margin=19mm]{geometry}
\usepackage{amsmath,amssymb,mathtools}
\usepackage[scr=rsfs,cal=euler]{mathalfa}
\usepackage{xfrac}
\usepackage[unicode,hidelinks,bookmarks=false]{hyperref}
\newcommand{\ie}{i.e.\ }
\newcommand{\cf}{cf.\ }
\newcommand{\resp}{respectively\ }
\newcommand{\Subsection}{\subsection}
\providecommand{\nobreakdash}{\nobreak\hbox{-}}
\setlength{\emergencystretch}{2em}
\multlinegap0pt
\usepackage{lmodern}
\usepackage{enumitem}
\usepackage{cleveref}
\newcommand{\crefrangeconjunction}{--}
\usepackage{tikz-cd}
\numberwithin{equation}{section}
\newtheorem{theorem}{Theorem}[section]
\newtheorem{lemma}[theorem]{Lemma}
\newtheorem{claim}[theorem]{Claim}
\newtheorem{proposition}[theorem]{Proposition}
\newtheorem{corollary}[theorem]{Corollary}
\newtheorem*{theorem*}{Theorem}
\newtheorem*{question}{Question}
\newtheorem*{problem}{Problem}
\newtheorem{definition}[theorem]{Definition}
\newtheorem{Remark}[theorem]{Remark}
\newtheorem{example}[theorem]{Example}
\newtheorem*{conjecture}{Conjecture}
\newcommand{\dl}{\operatorname{dl}}
\newcommand{\Fit}{\operatorname{F}}
\newcommand{\Fith}{\operatorname{Fh}}
\newcommand{\langleG}{\langle\,\cdot\,,\,\cdot\,\rangle}
\newcommand{\N}{\mathrm{N}}
\newcommand{\C}{\mathrm{C}}
\newcommand{\frattini}{\Phi}
\newcommand{\soc}{\mathrm{Soc}}
\newcommand{\Soc}{\mathrm{Soc}}
\newcommand{\Endd}{\mathrm{End}}
\newcommand{\End}{\mathrm{End}}
\newcommand{\Z}{\mathbb{Z}}
\newcommand{\cc}{\mathrm{cc}}
\newcommand{\gam}{\gamma}
\newcommand{\gen}[1]{\langle #1 \rangle}
\newcommand{\normcl}[1]{\langle #1^G \rangle}
\newcommand{\rls}{R_{\textup{L}\mathfrak{S}}}
\newcommand{\rad}{\mathcal{S}}
\newcommand{\AS}{\mathfrak{AS}}
\newcommand{\Mon}{\text{Mon}}
\newcommand{\Alt}{\operatorname{Alt}}
\newcommand{\Sym}{\operatorname{Sym}}
\newcommand{\FS}{\mathcal{FS}}
\newcommand{\Rls}{R_{\LS}}
\newcommand{\HP}{\mathrm{HP}}
\newcommand{\LS}{\textup{L}\mathfrak{S}}
\DeclareMathOperator{\Aut}{Aut}
\DeclareMathOperator{\Ker}{Ker}
\numberwithin{equation}{section}
\begin{document}
\title[Locally Solvable Radicals via Wilson Radical Sets and Subgro]{Locally Solvable Radicals via Wilson Radical Sets and Subgroup Lattices}
\author{Cao Minh Nam}
\date{}
\maketitle
\noindent\emph{Source and licence.} Locally Solvable Radicals via Wilson Radical Sets and Subgroup Lattices. Version arXiv:2608.04500v1 (CC BY 4.0). This material is distributed under the Creative Commons Attribution 4.0 International licence, http://creativecommons.org/licenses/by/4.0/, as recorded in the arXiv OAI licence metadata for this exact version and shown on the abstract page https://arxiv.org/abs/2608.04500v1. Full rights notice: licenses/arxiv-2608.04500-CC-BY-4.0.txt.\par\smallskip
\noindent\textit{Editorial scope.} Counted unit: the original \texttt{theorem} environment \texttt{theo:locally-finite} at source lines 249--260 together with its complete proof at lines 262--288; the delivered document keeps source lines 149--289 so that every referenced label and every citation resolves inside it. Native editable source is the paper's own concatenated TeX; the AMS wrapper is editorial formatting only. No publisher class, upstream script or unrelated artwork is redistributed.
\section{Locally (solvable-by-finite) groups}

In this section, we apply Wilson's characterization of the solvable radical in finite groups to locally (solvable-by-finite) groups.  
For each Wilson sequence $\omega$, we define a set $\mathcal W_\omega(G)$ through the eventual vanishing of the corresponding word values. 
Although this set is a priori dependent on $\omega$, we prove that, for every locally (solvable-by-finite) group $G$, 
$$ R_{\mathrm{L}\mathfrak S}(G) = R_{\mathrm{L}\mathfrak R}(G) = S(G) = \mathcal W_\omega(G).$$ 
Thus, the locally solvable radical is characterized by the condition defining $S(G)$. 
Analogous characterizations in terms of conjugates and commutators are established later in this section. 

\subsection{Wilson radical sets and local radicals} \hfill

\medskip
We begin by recalling the finite result on which the argument rests. 
Let $F(x,y)$ be the free group on two generators $x$ and $y$. 
A sequence $(w_n)$ in $F(x,y)$ is \emph{profinitely convergent} if, for every finite group $E$, there is a positive integer $N$ such~that $w_i(a,b)=w_j(a,b)$ for all $a,b\in E$ and all $i,j\geq N$. 

\begin{proposition}[{\cite[Eq.~(2) and Theorem]{Wilson2011}}]
\label{prop:Wilson-sequence}
There exists a sequence of~words $(w_n(x,y))_{n\geq 1}$ in the free group $F(x,y)$ with the following properties:
\begin{enumerate}[label=\textup{(\arabic*)}]
    \item The sequence $(w_n)_{n\geq 1}$ is profinitely convergent.
    \item For every integer $m\geq 1$, there exists an integer $N=N(m)$ such that
          $$w_n\in \left\langle [x,y]^{F(x,y)}\right\rangle^{(m)}$$ 
          for all $n\geq N$.
    \item For every finite group $G$ and every $g\in G$, one has $g$ belongs to the solvable radical $R(G)$ if and only if, for each $h\in G$, there exists an integer
          $N=N(g,h)$ such~that $w_n(g,h)=1$ for all $n\geq N$.
\end{enumerate}
\end{proposition}
  
Let $\Omega$ be the set of all sequences of words satisfying conditions \textup{(1)}--\textup{(3)} of Proposition~\ref{prop:Wilson-sequence}.  
For $\omega=(w_n(x,y))_{n\geq 1}\in\Omega$ and for an arbitrary group $G$, define the \emph{Wilson radical set associated with $\omega$}, henceforth
called simply the \emph{Wilson radical set}, by 
$$\mathcal W_\omega(G)=\left\{\,g\in G \mid \forall h\in G,\ \exists N(g,h):\ \forall n\geq N(g,h),\ w_n(g,h)=1 \,\right\}.$$

Wilson also extended Proposition~\ref{prop:Wilson-sequence}(3) to linear groups, drawing on Gruenberg's theorem \cite[Theorem~2]{Gruenberg66} that a finitely generated linear group is solvable if and only if all of its finite quotients are solvable.  We note this consequence for later use.

\begin{theorem}
\label{theo:Wilsonradsetlineargroup}
Let $\mathbb F$ be a field and let $G\leq \mathrm{GL}_n(\mathbb F)$.  Then, for every $\omega\in\Omega$,
$$\mathcal W_\omega(G)=R(G).$$
\end{theorem}

We now state two elementary facts about Wilson radical sets.

\begin{lemma}\label{lem:SsubseteqW}
Let $G$ be a group, and let $\omega\in\Omega.$ Then the following properties~hold:
\begin{enumerate}[label=\textup{(\arabic*)}]
    \item $S(G) \subseteq \mathcal{W}_{\omega}(G)$.
    \item The subgroup generated by $ \mathcal{W}_{\omega}(G)$ is characteristic in $G$.
\end{enumerate}
\end{lemma}
    
\begin{proof} (1) 
Fix $g\in S(G)$, let $h\in G$ be arbitrary, and put $H=\langle g,h\rangle$. 
Since $[g,h]\in\langle g^H\rangle$ and $\langle g^H\rangle\trianglelefteq H$, it follows that
  $$\bigl\langle [g,h]^H\bigr\rangle
  \leq
  \langle g^H\rangle.$$
Moreover, the solvability of $H$ implies that the closure $\bigl\langle [g,h]^H \bigr\rangle$ has derived length $m \geq 0$. 
Write $\omega=(w_n(x,y))_{n\geq 1}$. 
Since   $\bigl\langle [g,h]^H\bigr\rangle^{(m)}=1,$ Proposition~\ref{prop:Wilson-sequence}\textup{(2)} provides an integer $N=N(m)$ such that $w_n(g,h)=1$ for every $n\geq N$. 
This forces $ g \in \mathcal{W}_{\omega}(G) $. 
Consequently  $S(G) \subseteq \mathcal{W}_{\omega}(G).$

    \medskip

        (2) Fix an automorphism $ \varphi $ of $ G $. 
        It suffices to verify the inclusion $$\varphi\bigl(\mathcal{W}_{\omega}(G)\bigr) \subseteq \mathcal{W}_{\omega}(G).$$
 For any $a \in G$, the surjectivity of $\varphi$ allows us to write $a = \varphi(h)$ for some $h \in G$. 
 For $g\in\mathcal W_\omega(G)$, there exists an integer $N=N(g,h)$ such~that
  $$w_n(g,h)=1$$ whenever $n\geq N$.
By functoriality of word values, for every $n\geq N$, we~have $$ w_n(\varphi(g),a) = w_n(\varphi(g),\varphi(h)) = \varphi\bigl(w_n(g,h)\bigr) = 1. $$  
Therefore $\varphi(g)\in\mathcal W_\omega(G)$. The proof is now complete.
\end{proof}

\begin{Remark} \label{rm:24}
It follows from Theorem~\ref{theo:Wilsonradsetlineargroup} that, for every linear group $G$, the set $\mathcal W_\omega(G)$ is independent of the choice of $\omega\in\Omega$.
\end{Remark}

An \emph{ascending series} of a group $G$ is a family of subgroups $(G_\alpha)_{\alpha\leq\gamma}$, indexed by an ordinal $\gamma$, satisfying the following four conditions:
\begin{enumerate}
\item[(i)] $G_0=1$ and $G_\gamma=G$;
\item[(ii)] $G_\alpha\leq G_\beta$ whenever $\alpha\leq\beta\leq\gamma$;
\item[(iii)] $G_\alpha\trianglelefteq G_{\alpha+1}$ for every $\alpha<\gamma$;
\item[(iv)] $G_\lambda=\bigcup_{\alpha<\lambda}G_\alpha$ for every limit ordinal $\lambda\leq\gamma$.
\end{enumerate}
The subgroups $G_\alpha$ are
called the \emph{terms} of the series, and the quotient groups
$G_{\alpha+1}/G_\alpha,
\,
\alpha<\gamma,$
are called its \emph{factors}. A subgroup $A\leq G$ is \emph{ascendant} in $G$ if it occurs as a term of some ascending series of $G$.

Recall that a group is \emph{radical} if it has an ascending series with  locally nilpotent factors. For locally (solvable-by-finite) groups, local radicality and local solvability coincide by \cite[Theorem~2.6]{CNH2025}.

\begin{lemma} \label{lem:locradislocsolonlocsolbyfinite} Let $G$ be a locally \textup{(}solvable-by-finite\textup{)} group. Then $G$ is locally radical if and only if it is locally solvable. 
\end{lemma} 

We can now state the main identification for this class. 


\begin{theorem} \label{theo:locally-finite} 
Every locally \textup{(}solvable-by-finite\textup{)} group $G$ admits
the locally solvable radical.
Moreover, for each $\omega\in\Omega$,
  $$R_{\mathrm{L}\mathfrak{S}}(G)
  =
  R_{\mathrm{L}\mathfrak{R}}(G)
  =
  S(G)
  =
  \mathcal{W}_{\omega}(G).$$
\end{theorem}

\begin{proof}
By \cite[Lemma~4.6]{CNH2025}, the locally solvable radical $R_{\LS}(G)$ exists as the subgroup generated by all locally solvable normal subgroups of $G$. In view of Lemma~\ref{lem:locradislocsolonlocsolbyfinite}, it coincides with the $\textup{L}\mathfrak{R}$-radical $R_{\textup{L}\mathfrak{R}}(G)$. 

Next we claim that $R_{\LS}(G) \subseteq S(G)$. Indeed,
let $g \in R_{\LS}(G)$ and $x \in G$.  Observe that
$\langle g,x\rangle
\big/
\bigl(\langle g,x\rangle\cap R_{\LS}(G)\bigr)$
is cyclic. According to \cite[Lemma 2.5]{CNH2025}, the intersection $\langle g, x \rangle \cap R_{\LS}(G)$ is solvable.  This forces  the subgroup $ \langle g, x \rangle $ to be solvable, and hence the inclusion $ R_{\LS}(G) \subseteq S(G) $ holds.   Combining this inclusion with
Lemma~\ref{lem:SsubseteqW}, we obtain
$R_{\LS}(G)
\subseteq
S(G)
\subseteq
\mathcal{W}_{\omega}(G).$

Therefore, to complete the proof, it remains to verify $\mathcal{W}_{\omega}(G) \subseteq R_{\LS}(G).$
Put
$M=\bigl\langle \mathcal W_\omega(G)\bigr\rangle.$
We shall show that $M$ is locally solvable. Consider elements $g_1,\ldots,g_m\in\mathcal W_\omega(G)$ and put
$L=\langle g_1,\ldots,g_m\rangle.$ Since $G$ is locally (solvable-by-finite), the group $L$ is solvable-by-finite. Hence there exists a solvable normal subgroup $A\trianglelefteq L$ such that $Q=L/A$ is finite. Let $\pi\colon L\twoheadrightarrow Q$ be the canonical epimorphism. We claim that $\pi(g_i)\in R(Q)$ for every $i=1,\ldots,m$. Fix $i$ and let $q\in Q$. Choose $\ell\in L$ such that $\pi(\ell)=q$. Since $g_i\in\mathcal W_\omega(G)$, the definition of the Wilson radical set provides an integer $N_i=N(g_i,\ell)$ such that $w_n(g_i,\ell)=1$ for every $n\geq N_i$. Therefore, $$w_n\bigl(\pi(g_i),q\bigr) = w_n\bigl(\pi(g_i),\pi(\ell)\bigr) = \pi\bigl(w_n(g_i,\ell)\bigr) = 1$$ for every $n\geq N_i$. 
Thus $\pi(g_i)\in\mathcal W_\omega(Q).$ 
On the other hand, since $Q$ is finite, it follows from Proposition~\ref{prop:Wilson-sequence}\textup{(3)} that $\pi(g_i)\in R(Q).$

Now $Q$ is generated by $\pi(g_1),\ldots,\pi(g_m)$, all of which belong to $R(Q)$. Hence $Q=R(Q),$ so $Q$ is solvable. Since $A$ is solvable, it follows that $L$ is solvable. Therefore $M$ is locally solvable. Moreover, 
$M$ is characteristic in $G$ by Lemma~\ref{lem:SsubseteqW}. Thus $M$ is a normal locally solvable subgroup of $G$, and so $M\leq R_{\LS}(G).$ In particular, $\mathcal W_\omega(G)\subseteq R_{\LS}(G),$ as required.
\end{proof}


\begin{thebibliography}{99}
\bibitem[CNH2025]{CNH2025} L.V. Chua, C.M. Nam, and B.X. Hai, \emph{On the solvable radical of groups}, J. Algebra \textbf{682} (2025), 787--803.
\bibitem[Gruenberg66]{Gruenberg66} K.W. Gruenberg, \emph{The Engel structure of linear groups}, J. Algebra \textbf{3} (1966), 291--303.
\bibitem[Wilson2011]{Wilson2011} J.S. Wilson, \emph{Characterization of the solvable radical by a sequence of words}, J. Algebra \textbf{326} (2011), no.~1, 286--289.
\end{thebibliography}
\end{document}
